\chapter{First-Order Logic Core}
\label{chap:fol-core}

This chapter describes the part of the repository that builds first-order
syntax, derivations, semantics, and theory-level notions. The main source files
for this layer are:
\begin{itemize}
\item \texttt{Flypitch/FOL/Syntax.lean}
\item \texttt{Flypitch/FOL/Formula.lean}
\item \texttt{Flypitch/FOL/Proof.lean}
\item \texttt{Flypitch/FOL/Semantics.lean}
\item \texttt{Flypitch/FOL/Theory.lean}
\item \texttt{Flypitch/FOL/Bounded.lean}
\end{itemize}

\section{Languages, Terms, And Formulas}

\begin{definition}[First-Order Language]
A first-order language consists of families of function and relation symbols,
indexed by arity. In Lean this is the structure \texttt{Language}, with fields
\texttt{functions} and \texttt{relations}.
\end{definition}

The syntax uses de Bruijn variables. This means that bound variables are
represented by natural numbers rather than by names. The advantage is that
substitution and renaming become explicit structural operations, which is
exactly what later completeness and Henkin arguments need.

There is one unusual design choice that is worth explaining early. Terms and
formulas are built in \emph{partially applied} form. A term of type
\texttt{preterm L l} is a term in the language $L$ that still expects $l$
arguments before it becomes closed. Likewise \texttt{preformula L l} is a
formula that still expects $l$ relation arguments. Closed terms and formulas
are obtained at arity $0$:
\[
\texttt{term L := preterm L 0},
\qquad
\texttt{formula L := preformula L 0}.
\]

This packaging lets the development treat symbols uniformly and avoids a large
amount of repeated arity bookkeeping.

\begin{definition}[Structural Operations]
The core structural operations are:
\begin{itemize}
\item \texttt{lift\_term\_at} and \texttt{lift\_formula\_at}, which raise free
variables above a cutoff;
\item \texttt{subst\_term} and \texttt{subst\_formula}, which substitute a term
for a chosen variable;
\item a large collection of commutation lemmas showing how lifting and
substitution interact.
\end{itemize}
These lemmas are the bookkeeping layer needed for semantics, weakening, and
eventually Henkin witness constructions.
\end{definition}

\section{Natural Deduction}

\begin{definition}[Derivability]
For a set $\Gamma$ of formulas, the inductive type \texttt{prf} formalizes
natural-deduction derivations of $A$ from $\Gamma$. Its rules
include:
\begin{itemize}
\item assumption,
\item implication introduction and elimination,
\item ex falso,
\item universal introduction and elimination,
\item reflexivity of equality,
\item substitution of equals into formulas.
\end{itemize}
The notation $\Gamma \vdash A$ abbreviates a term of type
\texttt{prf} with assumptions $\Gamma$ and conclusion $A$, and
$\Gamma \vdash' A$ is the proposition that such a derivation exists.
\end{definition}

The file \texttt{Flypitch/FOL/Proof.lean} then develops the standard structural
machinery around this inductive system. In particular it proves weakening for
derivations and packages familiar derived rules such as conjunction and
disjunction introduction and elimination. This is the proof-theoretic core used
throughout the compactness and completion arguments.

\section{Structures And Satisfaction}

\begin{definition}[Structure]
An \texttt{L}-structure consists of a carrier type together with
interpretations of all function and relation symbols of the language. The Lean
name is \texttt{Structure L}.
\end{definition}

Given a structure $M$ and a valuation $v : \mathbb{N} \to M$, the development
defines:
\begin{itemize}
\item \texttt{realize\_term v t}, the interpretation of a term,
\item \texttt{realize\_formula v f}, the truth value of a formula,
\item \texttt{satisfied\_in M f}, satisfaction of a closed formula in a
structure,
\item $\Gamma \models A$, semantic consequence.
\end{itemize}

The main semantic bookkeeping theorems identify syntactic substitution with
semantic reassignment of valuations. This is the standard bridge that lets the
natural-deduction rules line up with model-theoretic semantics.

\begin{proposition}[Soundness]
\label{prop:formula-soundness}
If a formula $A$ is derivable from a set of assumptions $\Gamma$, then $A$ is a
semantic consequence of $\Gamma$.
\[
\Gamma \vdash A \implies \Gamma \models A.
\]
\end{proposition}

\begin{proof}
This is the theorem \texttt{formula\_soundness}. The proof is by induction on
the derivation. Each natural-deduction rule is checked directly against the
semantic definition of truth for formulas. The substitution lemmas proved
earlier are exactly what make the quantifier and equality cases go through.
\end{proof}

\section{Sentences, Theories, And Bounded Syntax}

Later chapters work primarily with closed formulas. For that reason
\texttt{Flypitch/FOL/Theory.lean} introduces:
\begin{itemize}
\item \texttt{sentence L}, the subtype of formulas whose free variables are all
below $0$,
\item \texttt{Theory L}, a set of sentences,
\item theory-level provability and satisfaction,
\item the predicates \texttt{is\_consistent} and \texttt{is\_complete}.
\end{itemize}

This file also isolates the notion of a term or formula being bounded by a
given number of free variables. The companion file
\texttt{Flypitch/FOL/Bounded.lean} packages bounded terms and formulas as
subtypes. These wrappers are not cosmetic. They are the later input to the
Henkin construction, where one must talk uniformly about formulas with a fixed
number of open variables and then substitute closed witnesses into them.

\begin{definition}[Bounded Syntax]
For each natural number $n$, the development defines:
\begin{itemize}
\item \texttt{bounded\_term L n},
\item \texttt{bounded\_formula L n},
\item \texttt{closed\_term L := bounded\_term L 0}.
\end{itemize}
These types package syntax together with proofs that all free variables lie
below the specified bound.
\end{definition}

\section{Sanity Check: Abelian Groups}

The file \texttt{Flypitch/Examples/Abel.lean} provides a compact regression
target for the logic core. It defines a language with a constant symbol and a
binary function symbol, writes down the usual axioms for abelian groups, and
interprets them in the integers.

\begin{proposition}[Integers Satisfy The Abelian-Group Axioms]
The structure \texttt{intStructure} satisfies the theory \texttt{TAb} of the
chosen abelian-group language.
\end{proposition}

\begin{proof}
The file proves the axioms one by one:
\begin{itemize}
\item \texttt{int\_satisfies\_assoc}
\item \texttt{int\_satisfies\_zeroRight}
\item \texttt{int\_satisfies\_zeroLeft}
\item \texttt{int\_satisfies\_inv}
\item \texttt{int\_satisfies\_comm}
\end{itemize}
It then bundles them into \texttt{int\_satisfies\_TAb}. Combining this with
Proposition~\ref{prop:formula-soundness} yields
\texttt{int\_satisfies\_of\_prf}: every theorem of \texttt{TAb} is true in the
integer model.
\end{proof}

This example plays two roles. It checks that the proof and semantics layers
interact correctly, and it gives a concrete model-theoretic example before the
blueprint moves on to the more abstract compactness and completion machinery.
